% Faithful editable transcription of official CCC2024 article32, pp.32:5--32:8.
\section*{Author circuit and hashing definitions}
\paragraph{Definition 8 (AC circuit).} AC circuit is a circuit with input gates, NOT gates, unbounded fan-in AND/OR gates, and (possibly more than one) output gates. The size of a circuit is defined by the number of AND/OR gates. The depth of a circuit is defined by the largest number of AND/OR gates on any path from an input gate to an output gate.
\paragraph{Definition 9 (L-uniform AC circuit family).} For functions $S,d:\mathbb Z^+\to\mathbb R^+$, we say a collection of circuits $\{C_n\}_{n\in\mathbb Z^+}$ is an $S$-size $d$-depth L-uniform AC circuit family, if each $C_n$ has size $\le S(n)$ and depth $\le d(n)$, and given the binary representation of $n$, the description of $C_n$ can be computed in uniform $O(\log n)$-space.
We need to mention that the number of input gates in $C_n$ is not necessarily $n$. Also note that since we can encode a tuple of $O(1)$ many integers to a single integer, we can also consider circuit collections with a tuple of integers as an index.
\paragraph{Definition 10 (Complexity Class L-AC$^1$).} We say a language $L$ is in the class L-AC$^1$ if there exists a $\operatorname{poly}(n)$-size $O(\log n)$-depth L-uniform AC circuit family $\{C_n\}$ such that $C_n$ computes $L$ on $n$-bit inputs.
\paragraph{Definition 11 (GapMaj).} For $n\in\mathbb Z^+$ and $a,b\in\mathbb R$ such that $0\le a<b\le n$, define the promise problem $\operatorname{GapMaj}[a,b]$ on $n$ bits as follow:
\[\operatorname{GapMaj}[a,b](x_1,\ldots,x_n):=\begin{cases}\mathrm{YES}&\text{if }x_1,\ldots,x_n\text{ contains }\ge b\text{ 1's}\\\mathrm{NO}&\text{if }x_1,\ldots,x_n\text{ contains }\le a\text{ 1's}\\\bot&\text{otherwise.}\end{cases}\]
We will use pairwise independent hash function as a tool for approximate counting in AC. We shall use the following construction based on convolution, which was also used in [22].
\paragraph{Definition 12 (Convolution-Based Pairwise Independent Hash Function).} Suppose $m$ is a power of 2. Define $H_m:[m^3]\times[m]\to[m]$ by: for $(k,x)\in[m^3]\times[m]$, let $x_1\cdots x_{\log m}$ be the binary representation of $x-1$, let $a_1\cdots a_{2\log m}b_1\cdots b_{\log m}$ be the binary representation of $k-1$, let $y_j:=(\sum_{i=1}^{\log m}a_{i+j}x_i+b_j)\bmod2$ for $j\in[\log m]$, then define $H_m(k,x)$ by letting $y_1\cdots y_{\log m}$ be the binary representation of $H_m(k,x)-1$.
\paragraph{Theorem 13.} $H_m$ is Pairwise Independent Hash Function in the following sense: for any $1\le i<j\le m$, when $k$ is sampled from the uniform distribution over $[m^3]$, the joint distribution of $(H_m(k,i),H_m(k,j))$ is identical to the uniform distribution over $[m]\times[m]$.
% STATEMENT_BEGIN
\paragraph{Theorem 14.} Let $n,a,b\in\mathbb N$ such that $0\le a<b\le n$. Then there exists a $\operatorname{poly}(n)$-size $O\left(\frac{\log\frac b{b-a}}{\log\log n}+1\right)$-depth L-uniform AC circuit family $\{C_{n,a,b}\}$ that computes $\operatorname{GapMaj}[a,b]$ on $n$ bits.
% STATEMENT_END
% PROOF_BEGIN
The proof depends on the next few Lemmas.
\paragraph{Lemma 15 ([2]).} Let $n\in\mathbb Z^+$. Then there exists $\operatorname{poly}(n)$-size $O(1)$-depth L-uniform AC circuit family $\{C_n^{(0)}\}$ that computes $\operatorname{GapMaj}[n/3,2n/3]$ on $n$ bits.
\paragraph{Lemma 16 (Exact Counting).} Let $n,\ell\in\mathbb Z^+$ such that $n\ge\ell$. Then there exists a $\operatorname{poly}(n)$-size $O\left(\frac{\log\ell}{\log\log n}+1\right)$-depth L-uniform AC circuit family $\{E_{n,\ell}\}$ such that on $\ell$ bits of input, $E_{n,\ell}$ outputs the exact number of 1's over the input bits, in binary form.
We remark that in Lemma 16, $n$ is only used to bound the size of the circuit. Also, $n,\ell$ are not necessarily polynomially related.
\begin{proof}
We only need to show how to compute sum of $O(\sqrt{\log n})$ many $O(\log n)$-bit\footnote{Here ``$O(\log n)$-bit integers'' refers to integers which has $O(\log n)$-bits in its binary representation.} non-negative integers in $O(1)$-depth, then by divide-and-conquer we can compute sum of $\ell$ bits in $O\left(\frac{\log\ell}{\log\log n}+1\right)$-depth.
View the $O(\log n)$-bit integers as $2^{\lfloor\sqrt{\log n}\rfloor}$-base $O(\sqrt{\log n})$-digit integers. We use the grade-school algorithm to sum $O(\sqrt{\log n})$ integers as follow. We first guess the result and all carry-bits, which involve at most $O(\sqrt{\log n})\cdot O\left(\log\left(\sqrt{\log n}\cdot2^{\lfloor\sqrt{\log n}\rfloor}\right)\right)=O(\log n)$ bits, and thus has at most $\operatorname{poly}(n)$ choices. Then we can apply a local check on each digit, each local check involves at most $O(\log n)$ bits, and thus deciding whether all local checks are passed can be computed in $O(1)$-depth. Then we can take the result of the only guess that passes all local checks. The total cost is $O(1)$-depth.
\end{proof}
\paragraph{Lemma 17.} Let $n,a,b\in\mathbb N$ such that $0\le a<b\le n$. Then there exists a $\operatorname{poly}(n)$-size $O\left(\frac{\log\frac n{b-a}}{\log\log n}+1\right)$-depth L-uniform AC circuit family $\{C_{n,a,b}^{(1)}\}$ that computes $\operatorname{GapMaj}[a,b]$ on $n$ bits.
\begin{proof}
Only consider the case that $n$ is a power of 2, otherwise we can use a simple padding argument. By Lemma 15, it suffices to show how to reduce $\operatorname{GapMaj}[a,b]$ on $n$ bits to $\operatorname{GapMaj}[n^3/3,2n^3/3]$ on $n^3$ bits, via a $\operatorname{poly}(n)$-size $O\left(\frac{\log\frac n{b-a}}{\log\log n}+1\right)$-depth L-uniform AC circuit.
If $b-a\le4\sqrt n$ then we can directly compute the number of 1's exactly via Lemma 16. Below we only consider $b-a>4\sqrt n$.
Let $\ell:=\left\lceil\frac{12n^2}{(b-a)^2}\right\rceil$. Suppose the $\operatorname{GapMaj}[a,b]$ instance is $x_1,x_2,\ldots,x_n$. Let $H_n$ be the hash function defined in Definition 12. Define $y_1,\ldots,y_{n^3}$ as follow: for $i\in[n^3]$, let $y_i$ be 1 if at least $\frac{a+b}{2n}$ fraction of $x_{H_n(i,1)},\ldots,x_{H_n(i,\ell)}$ is 1, otherwise let $y_i$ be 0. Note that $y_1,\ldots,y_{n^3}$ can be computed via a $\operatorname{poly}(n)$-size $O\left(\frac{\log\ell}{\log\log n}+1\right)$-depth L-uniform AC circuit, by Lemma 16. Here $O\left(\frac{\log\ell}{\log\log n}+1\right)=O\left(\frac{\log\frac n{b-a}}{\log\log n}+1\right)$.
Let's do some simple calculations. Assume $p$ fraction of $x_1,\ldots,x_n$ is 1. Let $S_i$ be number of 1's in $x_{H_n(i,1)},\ldots,x_{H_n(i,\ell)}$. Then we have $\mathbb E_{i\sim[n^3]}[S_i]=p\ell$ and $\operatorname{Var}_{i\sim[n^3]}[S_i]\le\ell$. So if $p\le\frac an$, then
\[\Pr_{i\sim[n^3]}\left[S_i\ge\ell\cdot\frac{a+b}{2n}\right]\le\frac\ell{(\ell\cdot\frac{b-a}{2n})^2}=\frac{4n^2}{\ell(b-a)^2}\le\frac13.\]
Similarly if $p\ge\frac bn$ then
\[\Pr_{i\sim[n^3]}\left[S_i\le\ell\cdot\frac{a+b}{2n}\right]\le\frac13.\]
This means if $x_1,\ldots,x_n$ is YES/NO instance of $\operatorname{GapMaj}[a,b]$, then $y_1,\ldots,y_{n^3}$ is YES/NO instance of $\operatorname{GapMaj}[n^3/3,2n^3/3]$. The reduction is completed.
\end{proof}
\begin{proof}[Proof of Theorem 14]
We will try to reduce to Lemma 17. Suppose the $\operatorname{GapMaj}[a,b]$ instance is $x_1,x_2,\ldots,x_n$. We only consider the case $n$ is a power of 2, otherwise use a simple padding argument. We only consider the case $10\left(\frac b{b-a}\right)^2<\frac n{b-a}$ (or equivalently, $n(b-a)>10b^2$), otherwise we can directly apply Lemma 17.
Let $\ell:=\left\lfloor\frac{n(b-a)}{2b^2}\right\rfloor$. For $i\in[n^3]$, let $y_i:=x_{H_n(i,1)}\vee\cdots\vee x_{H_n(i,\ell)}$, here $H_n$ is the hash function defined in Definition 12. Then $y_1,\ldots,y_{n^3}$ can be computed via $\operatorname{poly}(n)$-size $O(1)$-depth L-uniform AC circuit.
Assume $p$ fraction of $x_1,\ldots,x_n$ is 1. Let $S_i$ be number of 1's in $x_{H_n(i,1)},\ldots,x_{H_n(i,\ell)}$. Then we have $\mathbb E_{i\sim[n^3]}[S_i]=p\ell$ and $\mathbb E_{i\sim[n^3]}[S_i^2]=\ell(\ell-1)p^2+\ell p\le\ell p+\ell^2 p^2$. Thus by
\[\frac{\mathbb E_{i\sim[n^3]}[S_i]^2}{\mathbb E_{i\sim[n^3]}[S_i^2]}\le\Pr_{i\sim[n^3]}[S_i\ge1]\le\mathbb E_{i\sim[n^3]}[S_i]\]
we know: if $p\le\frac an$, then $\Pr_{i\sim[n^3]}[S_i\ge1]\le\frac{\ell a}n$; if $p\ge\frac bn$, then
\[\Pr_{i\sim[n^3]}[S_i\ge1]\ge\frac{(\frac{\ell b}n)^2}{\frac{\ell b}n+(\frac{\ell b}n)^2}\ge\frac{\ell b}n-\left(\frac{\ell b}n\right)^2.\]
To summarize, if $x_1,\ldots,x_n$ is YES/NO instance of $\operatorname{GapMaj}[a,b]$, then $y_1,\ldots,y_{n^3}$ is YES/NO instance of
\[\operatorname{GapMaj}\left[\left\lfloor n^3\cdot\frac{\ell a}n\right\rfloor,\left\lceil n^3\cdot\left(\frac{\ell b}n-\left(\frac{\ell b}n\right)^2\right)\right\rceil\right].\]
Finally we observe that
\[\left(\frac{\ell b}n-\left(\frac{\ell b}n\right)^2\right)-\frac{\ell a}n=\ell\cdot\left(\frac{b-a}n-\frac{\ell b^2}{n^2}\right)\ge\frac{n(b-a)}{3b^2}\cdot\frac{b-a}{2n}=\frac{(b-a)^2}{6b^2}.\]
Thus by Lemma 17, $\operatorname{GapMaj}\left[\left\lfloor n^3\cdot\frac{\ell a}n\right\rfloor,\left\lceil n^3\cdot\left(\frac{\ell b}n-\left(\frac{\ell b}n\right)^2\right)\right\rceil\right]$ over $n^3$ bits can be computed via a $\operatorname{poly}(n)$-size $O\left(\frac{\log\frac b{b-a}}{\log\log n}+1\right)$-depth L-uniform AC circuit.
\end{proof}
% PROOF_END
